% Standalone analytic module. Requires the main paper's standard theorem environments.
% Bibliography additions are in analytic-calibration-bibliography.tex.
\providecommand{\norm}[2]{\lVert #1\rVert_{#2}}
\section{Analytic calibration of contact-weighted errors}
Let $\mu_m$ be uniform probability measure on $\{-1,1\}^m$. Write
\[
 f(x)=\prod_{i=1}^m x_i,\qquad
 \partial_i P(x)=\frac{P(x)-P(x^{\oplus i})}{2},\qquad
 \Delta=\sum_{i=1}^m\partial_i.
\]
Thus $\Delta\chi_S=|S|\chi_S$. All unqualified norms and expectations use $\mu_m$. For $1<t<2$, put
\[
 s_t=1+\frac{2}{\pi}\arcsin\!\left(\frac2t-1\right).
\]
We use the dimension-free scalar Bernstein inequality
\begin{equation}\label{ac:eq:bernstein}
 \norm{\Delta P}{t}\le C_t u^{s_t}\norm{P}{t}
 \qquad (\deg P\le u).
\end{equation}
This is part of the established theory of polynomial inequalities on the cube \cite{EI}. Section~\ref{ac:sec:log} supplies the needed version, including locally uniform dependence on $t$, from the imaginary-power estimate of Carbonaro and Dragi\v cevi\'c \cite{CD}.

Fix $0<\beta\le1/2$ and put $p=2-\beta$. A contact law is a joint distribution $\pi$ of $(U,V)$ in the cube whose first marginal is exactly $\mu_m$. Set $L=\dist(U,V)$. For a margin-separated function $F$, meaning $|F|\ge1$, define its parity-error set
\[
 E_F=\{x:f(x)F(x)<0\}.
\]
The contact-weighted error cost outside an exceptional set $B$ is
\begin{equation}\label{ac:eq:cost}
 \mathcal C_{\pi,B}(F)
 =\E_\pi\frac{\mathbf 1_{\{U\in E_F\setminus B\}}|F(U)|^p}{(L+1)^\beta}.
\end{equation}
No independence between the amplitudes, labels, and contact distance is assumed.

\begin{lemma}[Stable contact-cost calibration]\label{lem:stable-contact-calibration}
Fix $1<p_0<p$, $c_0>0$, $K<\infty$, and $C_*<\infty$. Suppose that:
\begin{enumerate}
\item $|F|\ge1$ and $\E[f\sign F]\ge c_0$;
\item $Q$ has degree at most $u\ge1$ and $\norm{F-Q}{p_0}\le e\le c_0/2$;
\item $\mu_m(B)\le\delta$ and
\begin{equation}\label{ac:eq:absorb}
 \delta^{1/p_0-1/2}(p_0-1)^{-u/2}\le\frac1{16};
\end{equation}
\item $\mathcal C_{\pi,B}(F)\le K$ and, for every $q\ge1$,
\begin{equation}\label{ac:eq:distance}
 \norm{L+1}{L^q(\pi)}\le C_*\bigl(\sqrt{mq}+q+1\bigr).
\end{equation}
\end{enumerate}
Then
\begin{equation}\label{ac:eq:theta}
 u\ge c\,m^{\theta(p_0)},\qquad
 \theta(p_0)=\frac{1-\beta/(2p)}{s_{p_0}},
\end{equation}
where $c>0$ depends only on the displayed fixed parameters.
\end{lemma}

The exceptional set may depend arbitrarily on $F,Q$ and the contact law. Its small probability is not sufficient on its own; condition~\eqref{ac:eq:absorb} is what controls its amplitude contribution. In particular, for $\delta=\rho^b$, it suffices to have $u\le a\log(1/\rho)$ with $a$ sufficiently small.

\begin{theorem}[Limiting exponent]\label{thm:endpoint-contact-calibration}
Under the hypotheses of Lemma~\ref{lem:stable-contact-calibration}, choose, for sufficiently large $m$,
$p_0=p-1/\log(em)$ and impose~\eqref{ac:eq:absorb} uniformly over these exponents. Then
\begin{equation}\label{ac:eq:endpoint}
 u\ge c\,\frac{m^{\theta_*}}{[\log(em)]^\kappa},\qquad
 \theta_* = \frac{1-\beta/(2p)}{s_p},\qquad
 \kappa=\frac{\beta}{2p s_p}.
\end{equation}
For $\beta=1/2$, these constants are
\[
 \theta_*=0.6851115714\ldots,\qquad
 \kappa=0.1370223142\ldots.
\]
\end{theorem}

\subsection{Proof of the calibration theorem}
We first record two simple consequences of parity.

\begin{lemma}[Parity and positivity]\label{ac:lem:parity}
If $|F|\ge1$, $g=\sign F$, and $\E[fg]\ge c_0$, then
\begin{equation}\label{ac:eq:influence}
 \E[g\Delta F]\ge I(g)\ge mc_0,
 \qquad I(g)=\sum_i\Pp\{g(X)\ne g(X^{\oplus i})\}.
\end{equation}
For every real $Q$ and $1\le t<\infty$,
\begin{equation}\label{ac:eq:positive}
 \norm{Q}{t}\le 2\norm{(fQ)_-}{t}+\frac2m\norm{\Delta Q}{t}.
\end{equation}
\end{lemma}
\begin{proof}
On an edge where $g$ changes sign, margin separation gives
$(g(x)-g(x^{\oplus i}))(F(x)-F(x^{\oplus i}))\ge4$.
Edge symmetrization proves the first inequality in~\eqref{ac:eq:influence}.
Also $\E[fg]=\E[f\partial_i g]$, so
$|\E[fg]|\le\E|\partial_i g|=\operatorname{Inf}_i(g)$ for every $i$.
Summing proves its second inequality; a square on the correlation is unnecessary.

Let $A$ be normalized cube adjacency and $q=fQ$. Since parity reverses on every edge,
$(I+A)q=2f\Delta Q/m$. Write its right side as $h$. Positivity of $A$ gives
$q_+\le A q_-+|h|$. Contractivity of $A$ and the triangle inequality give
$\norm{q}{t}\le\norm{q_+}{t}+\norm{q_-}{t}\le2\norm{q_-}{t}+\norm{h}{t}$.
\end{proof}

\begin{proof}[Proof of Lemma~\ref{lem:stable-contact-calibration}]
Set $X=\norm Q{p_0}$ and $W=\norm{F\mathbf 1_{E_F\setminus B}}{p_0}$.
H\"older's inequality with conjugate exponents $p/p_0$ and $p/(p-p_0)$ yields
\begin{align}
 W
 &\le K^{1/p}\norm{L+1}{L^q(\pi)}^{\beta/p},
 &q&=\frac{\beta p_0}{p-p_0}.\label{ac:eq:holder}
\end{align}
If $q<1$, its norm is bounded by the first moment. For fixed $p_0<p$, hypothesis~\eqref{ac:eq:distance} therefore gives
\begin{equation}\label{ac:eq:W}
 W\le C m^{\beta/(2p)}.
\end{equation}
The exact uniform first marginal is used in identifying the left side of~\eqref{ac:eq:holder}.

The usual degree-$u$ hypercontractive comparison gives
$\norm Q2\le(p_0-1)^{-u/2}X$. Consequently,
\[
 \norm{F\mathbf 1_B}{p_0}
 \le\delta^{1/p_0-1/2}\norm Q2+e
 \le X/16+e.
\]
Since the negative-part map is $1$-Lipschitz, Lemma~\ref{ac:lem:parity} implies
\begin{equation}\label{ac:eq:boot}
 X\le2W+X/8+4e+\frac2m\norm{\Delta Q}{p_0}.
\end{equation}
If $2C_{p_0}u^{s_{p_0}}/m\ge1/4$, then $u\ge c m^{1/s_{p_0}}$, which is stronger than the conclusion. Otherwise~\eqref{ac:eq:bernstein} absorbs the last term in~\eqref{ac:eq:boot}, giving $X\le C(W+e)$.

Each $\partial_i$ is an $L^1$ contraction, so
$\norm{\Delta(F-Q)}1\le me$. By~\eqref{ac:eq:influence},
\[
 \norm{\Delta Q}{p_0}\ge\E[g\Delta Q]\ge m(c_0-e)\ge mc_0/2.
\]
Combining this with~\eqref{ac:eq:bernstein},~\eqref{ac:eq:W}, and the bound on $X$ proves~\eqref{ac:eq:theta}.
\end{proof}

\begin{proof}[Proof of Theorem~\ref{thm:endpoint-contact-calibration}]
Now $q=\beta p_0\log(em)$ in~\eqref{ac:eq:holder}, and hence
\[
 W\le C m^{\beta/(2p)}[\log(em)]^{\beta/(2p)}.
\]
The Bernstein constants are bounded on a fixed compact interval of exponents containing all sufficiently large-$m$ values of $p_0$. Moreover,
$s_{p_0}-s_p=O(1/\log(em))$. We may assume $u\le m$, so
$u^{s_{p_0}}\le C u^{s_p}$. The preceding proof gives~\eqref{ac:eq:endpoint}.
\end{proof}


\section{A logarithmic Bernstein refinement}\label{ac:sec:log}
The next estimate makes the norm dependence in scalar Bernstein explicit. Its proof follows the fixed-vector interpolation method of Domelevo et al.\ \cite{Domelevo}, with strip harmonic measure in place of Gaussian damping. We make no priority claim for the interpolation refinement.

\begin{theorem}\label{ac:thm:log}
For $1<t<2$, set $\sigma_t=(2/\pi)\arcsin(2/t-1)$. Every nonzero real cube polynomial of degree at most $u\ge1$ satisfies
\begin{equation}\label{ac:eq:logbern}
 \norm{\Delta P}t\le C_t u
 \left[1+\log\frac{\norm P2}{\norm Pt}\right]^{\sigma_t}\norm Pt.
\end{equation}
The constants are bounded when $t$ ranges over a compact subinterval of $(1,2)$.
\end{theorem}
\begin{proof}
Let $0<\varepsilon\le(t-1)/t$ and define
\[
 q=\frac{t(1-\varepsilon)}{1-t\varepsilon/2},\qquad
 \frac1t=\frac\varepsilon2+\frac{1-\varepsilon}{q}.
\]
We first prove
\begin{equation}\label{ac:eq:mixed}
 \norm{\Delta P}t\le C_tu\varepsilon^{-\sigma_t}
 \norm P2^\varepsilon\norm Pq^{1-\varepsilon}.
\end{equation}
For $L=\Delta+\gamma I$, $\gamma>0$, the imaginary-power estimate of \cite[Proposition 11]{CD} is
\[
 \norm{L^{iv}}{q\to q}\le C(q^*-1)(1+|v|)^{1/2}
 e^{\phi_q|v|},\qquad \phi_q=\arcsin(2/q-1).
\]
Its constants are independent of the dimension and of $\gamma$. Fix a real $g$ with $\norm g{t'}=1$, put $N=1/\varepsilon$, and form the holomorphic bilinear pairing
\[
 G(z)=\E\left[(L^{N(1-z)}P)\,g_z\right],\qquad
 g_z=\sign(g)|g|^{t'((1-z)/2+z/q')}.
\]
The value is set to zero where $g=0$. At $\theta=1-\varepsilon$, it equals $\E[(LP)g]$. On the left boundary,
$|G(iv)|\le(u+\gamma)^N\norm P2$; on the right,
\[
 |G(1+iv)|\le C(q^*-1)(1+N|v|)^{1/2}
 e^{\phi_qN|v|}\norm Pq.
\]
On the finite cube $G$ is bounded on the closed strip, so the Poisson inequality applies to $\log|G|$.

The right-boundary harmonic-measure density at $\theta$ is
\[
 w_\varepsilon(v)=\frac{\sin(\pi\varepsilon)}{2(\cosh(\pi v)-\cos(\pi\varepsilon))},
 \qquad \int w_\varepsilon=1-\varepsilon.
\]
The elementary bounds needed here are
\begin{align*}
 \int |v|w_\varepsilon(v)\,dv
 &\le\frac{2\varepsilon}{\pi}\bigl(\log(1/\varepsilon)+1\bigr),\\
 \int\log(1+|v|/\varepsilon)w_\varepsilon(v)\,dv&\le C.
\end{align*}
For the first, expand
\[\frac{\sin a}{\cosh x-\cos a}=2\sum_{n\ge1}\sin(na)e^{-nx}\]
and integrate, obtaining
\[(2/\pi^2)\int_0^{\pi\varepsilon}-\log(2\sin(s/2))\,ds.\]
For the second, use $w_\varepsilon(v)\le\pi\varepsilon/[4(v^2+\varepsilon^2)]$ and rescale.

The left boundary has mass $\varepsilon$. Integrating the boundary logarithms consequently gives
\[
 |G(\theta)|\le C_t(u+\gamma)
 \norm P2^\varepsilon\norm Pq^{1-\varepsilon}
 \varepsilon^{-2\phi_q/\pi}.
\]
Here $q\ge1+(t-1)/2$, and $|\phi_q-\phi_t|\le C_t\varepsilon$, so the last power is at most $C_t\varepsilon^{-\sigma_t}$. Taking the dual supremum and then $\gamma\downarrow0$ proves~\eqref{ac:eq:mixed}.

Since $q<t$, $\norm Pq\le\norm Pt$. Take
$\varepsilon=\min\{(t-1)/t,[1+\log(\norm P2/\norm Pt)]^{-1}\}$ to obtain~\eqref{ac:eq:logbern}. All bounds used above are uniform on compact exponent intervals.
\end{proof}

Hypercontractivity gives $\log(\norm P2/\norm Pt)\le(u/2)\log(1/(t-1))$. Substitution recovers~\eqref{ac:eq:bernstein}, since $s_t=1+\sigma_t$. Theorem~\ref{ac:thm:log} therefore does not prove the general linear-degree Bernstein conjecture; it gives linear degree when the norm ratio is bounded.

